\documentclass[10pt,a4paper]{amsart}
\usepackage{amsmath,amssymb,amsthm,graphicx}
% Retained minimal preamble for the Stacks Project chapter more-algebra.tex (More on Algebra).
% The chapter file pinned-original-source.tex carries only \input{preamble} on its first line, so no
% preamble text lives inside the chapter itself.  The Stacks Project uses one single shared preamble for
% every chapter; the snapshot stages it identically as sources/stacks-chow/preamble.tex and
% sources/stacks-etale-cohomology/preamble.tex.  The block below is copied mechanically, line for line,
% from sources/stacks-chow/preamble.tex lines 14-18 (the Xy-pic setup) and 154-217 (the native theorem
% declarations together with the mathematical macros).  Nothing was added, removed or redefined; the
% wrapper class and the amsmath/amssymb/amsthm/graphicx base packages are supplied by the assembler.
% Chapter: more-algebra.tex -- More on Algebra.
% For commutative diagrams we use Xy-pic
\usepackage[all]{xy}

% We use 2cell for 2-commutative diagrams.
\xyoption{2cell}

% Theorem environments.
%
\theoremstyle{plain}
\newtheorem{theorem}[subsection]{Theorem}
\newtheorem{proposition}[subsection]{Proposition}
\newtheorem{lemma}[subsection]{Lemma}

\theoremstyle{definition}
\newtheorem{definition}[subsection]{Definition}
\newtheorem{example}[subsection]{Example}
\newtheorem{exercise}[subsection]{Exercise}
\newtheorem{situation}[subsection]{Situation}

\theoremstyle{remark}
\newtheorem{remark}[subsection]{Remark}
\newtheorem{remarks}[subsection]{Remarks}

\numberwithin{equation}{subsection}

% Macros
%
\def\lim{\mathop{\mathrm{lim}}\nolimits}
\def\colim{\mathop{\mathrm{colim}}\nolimits}
\def\Spec{\mathop{\mathrm{Spec}}}
\def\Hom{\mathop{\mathrm{Hom}}\nolimits}
\def\Ext{\mathop{\mathrm{Ext}}\nolimits}
\def\SheafHom{\mathop{\mathcal{H}\!\mathit{om}}\nolimits}
\def\SheafExt{\mathop{\mathcal{E}\!\mathit{xt}}\nolimits}
\def\Sch{\mathit{Sch}}
\def\Mor{\mathop{\mathrm{Mor}}\nolimits}
\def\Ob{\mathop{\mathrm{Ob}}\nolimits}
\def\Sh{\mathop{\mathit{Sh}}\nolimits}
\def\NL{\mathop{N\!L}\nolimits}
\def\CH{\mathop{\mathrm{CH}}\nolimits}
\def\proetale{{pro\text{-}\acute{e}tale}}
\def\etale{{\acute{e}tale}}
\def\QCoh{\mathit{QCoh}}
\def\Ker{\mathop{\mathrm{Ker}}}
\def\Im{\mathop{\mathrm{Im}}}
\def\Coker{\mathop{\mathrm{Coker}}}
\def\Coim{\mathop{\mathrm{Coim}}}

% Boxtimes
%
\DeclareMathSymbol{\boxtimes}{\mathbin}{AMSa}{"02}

%
% Macros for moduli stacks/spaces
%
\def\QCohstack{\mathcal{QC}\!\mathit{oh}}
\def\Cohstack{\mathcal{C}\!\mathit{oh}}
\def\Spacesstack{\mathcal{S}\!\mathit{paces}}
\def\Quotfunctor{\mathrm{Quot}}
\def\Hilbfunctor{\mathrm{Hilb}}
\def\Curvesstack{\mathcal{C}\!\mathit{urves}}
\def\Polarizedstack{\mathcal{P}\!\mathit{olarized}}
\def\Complexesstack{\mathcal{C}\!\mathit{omplexes}}
% \Pic is the operator that assigns to X its picard group, usage \Pic(X)
% \Picardstack_{X/B} denotes the Picard stack of X over B
% \Picardfunctor_{X/B} denotes the Picard functor of X over B
\def\Pic{\mathop{\mathrm{Pic}}\nolimits}
\def\Picardstack{\mathcal{P}\!\mathit{ic}}
\def\Picardfunctor{\mathrm{Pic}}
\def\Deformationcategory{\mathcal{D}\!\mathit{ef}}

% --- appended 2026-10-12 (Chief): the project own definitions of the reference/slogan/history
% display environments, copied mechanically from the shared Stacks preamble lines 8-12
% (\usepackage{verbatim} plus the three \newenvironment lines, which the chapter preamble keeps
% verbatim from the project preamble). Nothing added, removed or redefined: these environments
% carry no mathematical content and the project itself typesets them as comments. Without these
% three lines a counted statement whose own text uses them fails with Environment slogan undefined.
\usepackage{verbatim}
\newenvironment{reference}{\comment}{\endcomment}
\newenvironment{slogan}{\comment}{\endcomment}
\newenvironment{history}{\comment}{\endcomment}
\usepackage{cleveref}
\crefname{originallemma}{original lemma}{original lemmas}
\Crefname{originallemma}{Original lemma}{Original lemmas}
\makeatletter
\newlabel{algebra-lemma-construct-fp-module}{{05N5}{source}{}{sourcepointer-1}{}}
\newlabel{algebra-lemma-construct-fp-module@cref}{{[originallemma][0][]05N5}{[1][0][]source}{}{}{}}
\makeatother
\crefname{originallemma}{original lemma}{original lemmas}
\Crefname{originallemma}{Original lemma}{Original lemmas}
\makeatletter
\newlabel{lemma-det}{{36464}{source}{}{sourcepointer-2}{}}
\newlabel{lemma-det@cref}{{[originallemma][0][]36464}{[1][0][]source}{}{}{}}
\makeatother
\crefname{originallemma}{original lemma}{original lemmas}
\Crefname{originallemma}{Original lemma}{Original lemmas}
\makeatletter
\newlabel{lemma-flat-finite-type-valuation-ring-finite-presentation}{{5967}{source}{}{sourcepointer-3}{}}
\newlabel{lemma-flat-finite-type-valuation-ring-finite-presentation@cref}{{[originallemma][0][]5967}{[1][0][]source}{}{}{}}
\makeatother
\crefname{originallemma}{original lemma}{original lemmas}
\Crefname{originallemma}{Original lemma}{Original lemmas}
\makeatletter
\newlabel{lemma-perfect-over-regular-local-ring}{{20229}{source}{}{sourcepointer-4}{}}
\newlabel{lemma-perfect-over-regular-local-ring@cref}{{[originallemma][0][]20229}{[1][0][]source}{}{}{}}
\makeatother
\crefname{originallemma}{original lemma}{original lemmas}
\Crefname{originallemma}{Original lemma}{Original lemmas}
\makeatletter
\newlabel{lemma-perfect-to-K-group}{{36491}{source}{}{sourcepointer-5}{}}
\newlabel{lemma-perfect-to-K-group@cref}{{[originallemma][0][]36491}{[1][0][]source}{}{}{}}
\makeatother
\crefname{originallemma}{original lemma}{original lemmas}
\Crefname{originallemma}{Original lemma}{Original lemmas}
\makeatletter
\newlabel{lemma-structure-relatively-perfect}{{21980}{source}{}{sourcepointer-6}{}}
\newlabel{lemma-structure-relatively-perfect@cref}{{[originallemma][0][]21980}{[1][0][]source}{}{}{}}
\makeatother
\crefname{originallemma}{original lemma}{original lemmas}
\Crefname{originallemma}{Original lemma}{Original lemmas}
\makeatletter
\newlabel{lemma-valuation-ring-torsion-free-flat}{{4925}{source}{}{sourcepointer-7}{}}
\newlabel{lemma-valuation-ring-torsion-free-flat@cref}{{[originallemma][0][]4925}{[1][0][]source}{}{}{}}
\makeatother
\title{The Picard group of the generic fibre vanishes for a flat essentially finite type local algebra over a valuation ring with regular special fibre}
\author{The Stacks Project Authors}
\date{stacks-project:more-algebra.tex:local-snapshot-20261012; GNU FDL 1.2 or later}
\begin{document}
\maketitle
\setcounter{subsection}{1}
\noindent\textit{Editorial note.} All delivered mathematical text is the author's own, copied from the pinned Stacks Project chapter file more-algebra.tex; no statement or proof step was rewritten, reordered or removed.\par\medskip
\noindent\textit{Editorial note.} The chapter file pinned-original-source.tex carries only \textbackslash{}input\{preamble\} on its first line, so no preamble text lives inside the chapter. The retained file sources/stacks-more-algebra/retained-preamble.tex is a mechanical, line-for-line copy of the project's single shared preamble as staged in the snapshot; nothing was added, removed or redefined and no package was added.\par\medskip
\noindent\textit{Editorial note.} The source URL is the source-level Stacks Project URL, as required when more than one card is produced from one source; the per-claim tag is carried by claim\_key.\par\medskip
\noindent\textit{Editorial note.} The counted statement span contains the author's three enumerated hypotheses, delivered verbatim. The author uses italic markup for the word 'free' inside the counted proof span; that is the author's text and is delivered as it stands. The seven references inside the counted proof span are recorded as non-load-bearing navigation pointers.\par\medskip
\noindent\textit{Allowed standard prerequisites (editorial).}\begin{itemize}
\item Torsion free modules over a valuation ring are flat, and flat finite type modules are finitely presented. invoked at proof lines 36864-36868. Gives flatness and finite presentation of the descended module M.
\item Relative perfectness and perfection over a regular local ring. invoked at proof lines 36869-36874. Upgrades M to a perfect module with a finite free resolution.
\item The Euler characteristic on K\_0 and the determinant map. invoked at proof lines 36877-36881. Turn the finite free resolution into the class statement and then into triviality of L.
\end{itemize}\medskip
\ifdefined\hypertarget\hypertarget{sourcepointer-1}{}\fi
\noindent\textit{Source reference.} Original lemma 05N5: focused-seeds/AGP/sources/stacks-snapshot-20261011/tags.txt, source lines 7294--7294. Cross-chapter navigation pointer into the Stacks Project; the standard fact it supplies is applied inside the counted proof span and the argument resting on it stays there.\par\medskip
\ifdefined\hypertarget\hypertarget{sourcepointer-2}{}\fi
\noindent\textit{Source reference.} Original lemma 36464: sources/stacks-more-algebra/pinned-original-source.tex, source lines 36464--36464. Same-chapter navigation pointer; the statement it names is standard for the section and the argument that consumes it stays inside the counted proof span.\par\medskip
\ifdefined\hypertarget\hypertarget{sourcepointer-3}{}\fi
\noindent\textit{Source reference.} Original lemma 5967: sources/stacks-more-algebra/pinned-original-source.tex, source lines 5967--5967. Same-chapter navigation pointer; the statement it names is standard for the section and the argument that consumes it stays inside the counted proof span.\par\medskip
\ifdefined\hypertarget\hypertarget{sourcepointer-4}{}\fi
\noindent\textit{Source reference.} Original lemma 20229: sources/stacks-more-algebra/pinned-original-source.tex, source lines 20229--20229. Same-chapter navigation pointer; the statement it names is standard for the section and the argument that consumes it stays inside the counted proof span.\par\medskip
\ifdefined\hypertarget\hypertarget{sourcepointer-5}{}\fi
\noindent\textit{Source reference.} Original lemma 36491: sources/stacks-more-algebra/pinned-original-source.tex, source lines 36491--36491. Same-chapter navigation pointer; the statement it names is standard for the section and the argument that consumes it stays inside the counted proof span.\par\medskip
\ifdefined\hypertarget\hypertarget{sourcepointer-6}{}\fi
\noindent\textit{Source reference.} Original lemma 21980: sources/stacks-more-algebra/pinned-original-source.tex, source lines 21980--21980. Same-chapter navigation pointer; the statement it names is standard for the section and the argument that consumes it stays inside the counted proof span.\par\medskip
\ifdefined\hypertarget\hypertarget{sourcepointer-7}{}\fi
\noindent\textit{Source reference.} Original lemma 4925: sources/stacks-more-algebra/pinned-original-source.tex, source lines 4925--4925. Same-chapter navigation pointer; the statement it names is standard for the section and the argument that consumes it stays inside the counted proof span.\par\medskip
\begin{lemma}
\label{lemma-picard-group-generic-fibre-regular}
Let $R$ be a valuation ring with fraction field $K$
and residue field $\kappa$. Let $R \to A$ be a
homomorphism of rings such that
\begin{enumerate}
\item $A$ is local and $R \to A$ is local,
\item $A$ is flat and essentially of finite type over $R$,
\item $A \otimes_R \kappa$ regular.
\end{enumerate}
Then $\Pic(A \otimes_R K) = 0$.
\end{lemma}

\begin{proof}
Let $L$ be an invertible $A \otimes_R K$-module. In particular $L$ is
a finite module. There exists a finite $A$-module
$M$ such that $M \otimes_R K \cong L$, see
Algebra, Lemma \ref{algebra-lemma-construct-fp-module}.
We may assume $M$ is torsion free as an $R$-module.
Thus $M$ is flat as an $R$-module
(Lemma \ref{lemma-valuation-ring-torsion-free-flat}).
From Lemma \ref{lemma-flat-finite-type-valuation-ring-finite-presentation}
we deduce that $M$ is of finite presentation as an $A$-module
and $A$ is essentially of finite presentation as an $R$-algebra.
By Lemma \ref{lemma-structure-relatively-perfect}
we see that $M$ is perfect relative to $R$,
in particular $M$ is pseudo-coherent as an $A$-module.
By Lemma \ref{lemma-perfect-over-regular-local-ring}
we see that $M$ is perfect, hence
$M$ has a finite free resolution $F_\bullet$ over $A$.
It follows that $L$ is quasi-isomorphic to a finite complex
of {\it free} $A \otimes_R K$-modules. Hence by
Lemma \ref{lemma-perfect-to-K-group} we see that
$[L] = n[A \otimes_R K]$ in $K_0(A \otimes_R K)$
for some $n \in \mathbf{Z}$.
Applying the map of Lemma \ref{lemma-det}
we see that $L$ is trivial.
\end{proof}

\end{document}
